\input{../tex-inputs/latex-leseansicht-vorspann}

\section[Hugo von Hofmannsthal und Arthur Schnitzler an Gustav Schwarzkopf, 30. 6. {[}1902{]}]{L04428 Hugo von Hofmannsthal und Arthur Schnitzler an Gustav Schwarzkopf, 30. 6. {[}1902{]}}
\nopagebreak\mylabel{L04428v}
\rehead{ }\normalsize\beginnumbering\briefempfaengerindex{Schwarzkopf, Gustav@\textsc{Schwarzkopf, Gustav}!zzzSchnitzler, Arthur@\emph{von Arthur Schnitzler}!1902-06-302@{30. 6. {[}1902{]}}|(be}\briefempfaengerindex{Schwarzkopf, Gustav@\textsc{Schwarzkopf, Gustav}!zzzHofmannsthal, Hugo von@\emph{von Hugo von Hofmannsthal}!1902-06-302@{30. 6. {[}1902{]}}|(be}
\toendnotes[C]{\smallbreak\pagebreak[2]}
\correspDesc{Versand  durch Hugo von Hofmannsthal, Arthur Schnitzler am 30. 6. [1902] in Salzburg
\newline{}Übermittlung  am 30. [6. 1902] in Salzburg
\newline{}Erhalt  durch Gustav Schwarzkopf im Zeitraum [1. 7. 1902 – 5. 7. 1902?] in Wien}\toendnotes[C]{\smallbreak}
\Standort{DLA, A:Schnitzler, HS.1985.1.1897.}
\physDesc{Bildpostkarte, 80 Zeichen
\newline{}Handschrift Hugo von Hofmannsthal: blauer Buntstift, deutsche Kurrent
\newline{}Handschrift Arthur Schnitzler: Bleistift
\newline{}Versand: Stempel: »\nobreak{}\oindex{Hauptbahnhof Salzburg@\textbf{Hauptbahnhof Salzburg}|pwk}Salzburg
                                       Bahnh\textcolor{gray}{of}, 30 {[}6. 02{]}\nobreak{}«.  }\toendnotes[C]{\smallbreak}\pstart{}{\pb}\textsc{Herrn Gustav
                     Schwarzkopf}\pend{}\pstart{}\textsc{Wien}\oindex{Wien@\textbf{Wien}|pw}\pend{}\pstart{}\textsc{I Tiefer Graben
                     23}\oindex{Wien@\textbf{Wien}!I., Innere Stadt@\textbf{I., Innere Stadt}!Tiefer Graben 23@\textbf{Tiefer Graben 23}|pw}.\pend{}{\bigskip}
\pstart
           \centering{}{\pb}\textcolor{gray}{\textbf{Salzburg}}\oindex{Salzburg@\textbf{Salzburg}|pw}\pend
           \vspace{1em}
\pstart
           {\pb}Herzliche Grüße{\\[\baselineskip]}von \spacefill\mbox{Hugo}\pend
           \leftskip=0em{}
\pstart
           \label{K_L04428-1v}\edtext{30 VI}{\lemma{\textnormal{\emph{30 VI}}}\Cendnote{\textnormal{Die fehlende Jahresangabe ergibt sich
                     durch den gemeinsamen Aufenthalt in Salzburg\oindex{Salzburg@\textbf{Salzburg}|pwk} zwischen
                     28. 6. 1902 und 1. 7. 1902.}}}\label{K_L04428-1}\pend
           \selectlanguage{ngerman}\vspace{1em}\pstart \spacefill\mbox{{[}hs. Schnitzler:{]} Arthur}\pend{}\selectlanguage{ngerman}\endnumbering\briefempfaengerindex{Schwarzkopf, Gustav@\textsc{Schwarzkopf, Gustav}!zzzSchnitzler, Arthur@\emph{von Arthur Schnitzler}!1902-06-302@{30. 6. {[}1902{]}}|)be}\briefempfaengerindex{Schwarzkopf, Gustav@\textsc{Schwarzkopf, Gustav}!zzzHofmannsthal, Hugo von@\emph{von Hugo von Hofmannsthal}!1902-06-302@{30. 6. {[}1902{]}}|)be}\mylabel{L04428h}
\begin{anhang}
\end{anhang}\newcommand{\dateiname}{L04428}\newcommand{\titel}{Hugo von Hofmannsthal und Arthur Schnitzler an Gustav Schwarzkopf, 30. 6. [1902]}\newcommand{\editorInnen}{Herausgegeben von Jahnke, SelmaMüller, Martin Anton}\input{../tex-inputs/latex-leseansicht-abspann}
